\documentclass[12pt]{article}
\usepackage[a4paper, tmargin=3cm, lmargin=2.5cm, rmargin=2.5cm, bmargin=3cm]{geometry}
\usepackage{hyperref}

\usepackage[brazilian]{babel} % Remover linha em caso de escrita em inglês
\usepackage[fixlanguage]{babelbib} % Remover linha em caso de escrita em inglês


\hypersetup{
    colorlinks=true,
    linkcolor=black,
    filecolor=magenta,
    urlcolor=blue,
    citecolor=black,
}

\usepackage{natbib} % to cite references by surname and year
\usepackage{graphicx}
\usepackage{fancyhdr}
\pagestyle{fancy} % Ativar o estilo fancy
\fancyhf{} % Limpar cabeçalhos e rodapés
\renewcommand{\headrulewidth}{0pt} % Remover a linha do cabeçalho
\fancyfoot[C]{IC/UNICAMP \\ 2026 }


%%%%%%%%%%%%%%%%%%%%%%%%%%%%%%%%%%
% CONFIGURAÇÃO DA PRIMEIRA PÁGINA
%%%%%%%%%%%%%%%%%%%%%%%%%%%%%%%%%%


\begin{document}

\vspace{5mm}
\hrule

\vspace{3mm}

\begin{center}
    {\large\textbf{INSTITUTO DE COMPUTAÇÃO - UNICAMP}} \\
    \vspace{3mm}
    {\large\textbf{MO810A/MC959A - Sistemas Multiagentes com Modelos de Linguagem}} \\
    \vspace{3mm}
    {\large\textbf{Prof. Dr. Julio Cesar dos Reis}}
\end{center}

\vspace{3mm}
\hrule

\vspace{2.5cm}
\begin{flushright}
    \Large{\textbf{Fase 3 - Resultado Prático}} \\
    \vspace{1cm}
    \textit{
        $<$Título do projeto$>$ \\
    }
    \vspace{5cm}
    $<$Aluno(a)$>$ - $<$RA Aluno(a)$>$ \\
    $<$Aluno(b)$>$ - $<$RA Aluno(b)$>$
\end{flushright}

%%%%%%%%%%%%%%%%%%%%%%%%%%%%%%%%%%
%%%%%%%%%%%%%%%%%%%%%%%%%%%%%%%%%%


\newpage
\pagestyle{fancy} % 2. Ativar o estilo fancy
\fancyhf{} % 3. Limpar cabeçalhos e rodapés
\renewcommand{\headrulewidth}{0pt} % Remover a linha do cabeçalho

\fancyfoot[R]{\thepage}


%%%%%%%%%%%%%%%%%%%%%%%%%%%%%%%%%%
%       INÍCIO DAS SEÇÕES
%%%%%%%%%%%%%%%%%%%%%%%%%%%%%%%%%%

%%%%%%%%%%%%%%%%%%%%%%%%%%%%%%%%%%%%%%%%
\section{Tecnologias Utilizadas}
\label{sec:tecnologias}
%%%%%%%%%%%%%%%%%%%%%%%%%%%%%%%%%%%%%%%%

%isso deve ser removido
\textit{$<$Explique as tecnologias utilizadas na implementação do sistema multiagente: linguagem e bibliotecas, \textit{frameworks} de construção e orquestração de agentes, modelos de linguagem empregados (com provedor, versão e modo de acesso), bancos de dados e índices vetoriais, APIs e serviços externos expostos aos agentes como ferramentas, e a infraestrutura de execução. Indique as versões relevantes e justifique cada escolha em relação à metodologia definida na Fase 2.$>$}

%%%%%%%%%%%%%%%%%%%%%%%%%%%%%%%%%%%%%%%%
\section{Decisões Chaves de Implementação}
\label{sec:decisoes}
%%%%%%%%%%%%%%%%%%%%%%%%%%%%%%%%%%%%%%%%

%isso deve ser removido
\textit{$<$Explique o pipeline de desenvolvimento que foi conduzido, das fontes de dados até a saída entregue ao usuário. Explique as decisões chaves da implementação com as tecnologias da Seção \ref{sec:tecnologias}: como cada agente foi especificado (\textit{prompts}, papéis, escopo de decisão e critério de término), como a orquestração e a troca de mensagens foram realizadas, como as ferramentas foram definidas e como suas saídas são validadas, quais mecanismos de memória e de controle de execução, custo e falhas foram adotados. Explicite os ajustes em relação à metodologia da Fase 2, as alternativas descartadas e o motivo.$>$}

%%%%%%%%%%%%%%%%%%%%%%%%%%%%%%%%%%%%%%%%
\section{O Resultado Principal Obtido}
\label{sec:resultado}
%%%%%%%%%%%%%%%%%%%%%%%%%%%%%%%%%%%%%%%%

%isso deve ser removido
\textit{$<$Descreva o resultado da implementação. Explique o produto final: o que o sistema multiagente recebe como entrada, o que entrega como saída e qual caso de uso foi demonstrado \textit{fim a fim}. Apresente um diagrama arquitetural da solução, identificando agentes, ferramentas, fontes de dados e o fluxo de controle entre eles. Explique como executar sua solução: requisitos, instalação das dependências, configuração de credenciais e parâmetros, comando de execução e um exemplo de entrada com a saída correspondente. Indique também as limitações observadas na prática.$>$}

\subsection*{Código da Solução}
\label{sec:codigo}

%isso deve ser removido
\textit{$<$Informe abaixo os dados da \textit{release} entregue. Somente o código da \textit{release} com a \textit{tag} \texttt{fase3}, criada até o prazo de entrega, será avaliado; \textit{commits} posteriores serão ignorados. Confirme que o repositório é privado e que o professor e o monitor da disciplina foram adicionados como colaboradores. As instruções de execução descritas nesta seção devem corresponder ao \texttt{README} dessa \textit{release}.$>$}

\begin{itemize}
    \item \textbf{Link da \textit{release}:} $<$\url{https://github.com/usuario/repositorio/releases/tag/fase3}$>$
    \item \textbf{\textit{Tag}:} \texttt{fase3}
    \item \textbf{\textit{Hash} do \textit{commit}:} \texttt{$<$hash do commit da tag$>$}
\end{itemize}

%%%%%%%%%%%%%%%%%%%%%%%%%%%%%%%%%%%%%%%%
\section*{Apêndice}
\label{sec:anexo}
%%%%%%%%%%%%%%%%%%%%%%%%%%%%%%%%%%%%%%%%

\textit{$<$Use se necessário. Documentos e outras informações podem ser anexados para apoiar o entendimento da solução desenvolvida, como \textit{prompts} completos dos agentes, esquemas das ferramentas, exemplos de \textit{logs} de execução e trechos de código. $>$}



%%%%%%%%%%%%%%%%%%%%%%%%%%%%%%%%%%
%          BIBLIOGRAFIA
%%%%%%%%%%%%%%%%%%%%%%%%%%%%%%%%%%
\newpage
\bibliographystyle{aasjournal}
\bibliography{references}{}



\end{document}
